% GENERATED FILE. Edit canonical lessons, then run npm run generate:books.
% canonical-source-hash: fnv1a64:478f33bb0cb497a8
% canonical-lessons: HI-C555-phalna, HI-C555-dolna, HI-C555-tikana, HI-C555-giravat, HI-C555-patan

\chapter{Ups and Downs}
\label{ch:hi-ups-and-downs}
\hlchaptermodality{555}

\begin{chapteropening}
\textbf{By the end of this chapter:} \emph{I can talk about things rising, falling and holding steady.}

Complete the last lesson of chapter 555: I can talk about things rising, falling and holding steady.
\end{chapteropening}

\section[\texorpdfstring{\hi{फलना} (phalnā)}{फलना (phalnā)}]{\hi{फलना} (phalnā) — to bear fruit; to pay off}

\label{lesson:HI-C555-phalna}



\begin{quote}
Before the new one: say the Hindi for approximately, then the Hindi for actually.
\end{quote}

\subsection*{You'll want to know: \hi{फलना}}
\textbf{\hi{फलना}} — \emph{phalnā} — \textquotedblleft{}to bear fruit; to pay off\textquotedblright{}.

\begin{quote}
\hi{आँगन} \hi{का} \hi{पेड़} \hi{बहुत} \hi{फला।}
\end{quote}

\begin{quote}
\emph{Āṁgan kā peṛ bahut phalā.}
\end{quote}

\begin{quote}
The tree in the courtyard bore a lot of fruit.
\end{quote}

Past (it …): \hi{वह} \hi{पेड़} \hi{फला}, \emph{vah peṛ phalā}.

\begin{cousinweb}
\textbf{\hi{फलना}} (\emph{phalnā}) is a tadbhava, inherited from Sanskrit \emph{phalati}, \textquotedblleft{}bears fruit\textquotedblright{} (CDIAL 9057).
\end{cousinweb}

\begin{grammarlens}[title={What you've built}]
One more word for saying what happens if.
\end{grammarlens}

\subsection*{Your turn}
\begin{itemize}
\raggedright
  \item \emph{Say it:} \emph{phalnā}
  \item \emph{Say it:} \emph{phalnā}, once more
  \item \emph{Say it:} \emph{Āṁgan kā peṛ bahut phalā.} — The tree in the courtyard bore a lot of fruit.
  \item \emph{Recall:} say the Hindi for approximately, then the Hindi for actually, then say \emph{phalnā} again
\end{itemize}

\begin{tcolorbox}[breakable,colback=teal!4,colframe=teal!35!black,title={Before you move on}]
What does \hi{फलना} mean? (\textquotedblleft{}To bear fruit; to pay off\textquotedblright{}.) Say it once more.
\end{tcolorbox}
\section[\texorpdfstring{\hi{डोलना} (ḍolnā)}{डोलना (ḍolnā)}]{\hi{डोलना} (ḍolnā) — to sway; to rock; to waver}

\label{lesson:HI-C555-dolna}



\begin{quote}
Before the new one: say the Hindi for actually, then the Hindi for to bear fruit.
\end{quote}

\subsection*{You'll want to know: \hi{डोलना}}
\textbf{\hi{डोलना}} — \emph{ḍolnā} — \textquotedblleft{}to sway; to rock; to waver\textquotedblright{}.

\begin{quote}
\hi{तेज़} \hi{हवा} \hi{में} \hi{नाव} \hi{डोल} \hi{रही} \hi{थी।}
\end{quote}

\begin{quote}
\emph{Tez havā meṁ nāv ḍol rahī thī.}
\end{quote}

\begin{quote}
The boat was rocking in the strong wind.
\end{quote}

Past (it …): \hi{वह} \hi{झूला} \hi{डोला}, \emph{vah jhūlā ḍolā}.

\begin{cousinweb}
\textbf{\hi{डोलना}} (\emph{ḍolnā}) is a tadbhava, inherited from Sanskrit \emph{dōlāyatē}, \textquotedblleft{}moves to and fro\textquotedblright{} (CDIAL 6585). CDIAL gives the Hindi verb as \textquotedblleft{}to oscillate\textquotedblright{}.
\end{cousinweb}

\begin{grammarlens}[title={What you've built}]
One more word for saying what happens if.
\end{grammarlens}

\subsection*{Your turn}
\begin{itemize}
\raggedright
  \item \emph{Say it:} \emph{ḍolnā}
  \item \emph{Say it:} \emph{ḍolnā}, once more
  \item \emph{Say it:} \emph{Tez havā meṁ nāv ḍol rahī thī.} — The boat was rocking in the strong wind.
  \item \emph{Recall:} say the Hindi for actually, then the Hindi for to bear fruit, then say \emph{ḍolnā} again
\end{itemize}

\begin{tcolorbox}[breakable,colback=teal!4,colframe=teal!35!black,title={Before you move on}]
What does \hi{डोलना} mean? (\textquotedblleft{}To sway; to rock; to waver\textquotedblright{}.) Say it once more.
\end{tcolorbox}
\section[\texorpdfstring{\hi{टिकाना} (ṭikānā)}{टिकाना (ṭikānā)}]{\hi{टिकाना} (ṭikānā) — to rest (something) against; to steady}

\label{lesson:HI-C555-tikana}



\begin{quote}
Before the new one: say the Hindi for to bear fruit, then the Hindi for to sway.
\end{quote}

\subsection*{You'll want to know: \hi{टिकाना}}
\textbf{\hi{टिकाना}} — \emph{ṭikānā} — \textquotedblleft{}to rest (something) against; to steady\textquotedblright{}.

\begin{quote}
\hi{उसने} \hi{सीढ़ी} \hi{दीवार} \hi{से} \hi{टिका} \hi{दी।}
\end{quote}

\begin{quote}
\emph{Usne sīṛhī dīvār se ṭikā dī.}
\end{quote}

\begin{quote}
He leaned the ladder against the wall.
\end{quote}

Past, with \hi{ने} (he or she …): \hi{उसने} \hi{साइकिल} \hi{दीवार} \hi{से} \hi{टिकाई}, \emph{usne sāikil dīvār se ṭikāī}.

\begin{cousinweb}
\textbf{\hi{टिकाना}} (\emph{ṭikānā}) is the causative of \emph{ṭiknā}, \textquotedblleft{}to stop, remain\textquotedblright{}, which CDIAL files under \emph{*ṭakk}, \textquotedblleft{}remain, stop\textquotedblright{}, a reconstructed form (CDIAL 5420).
\end{cousinweb}

\begin{grammarlens}[title={What you've built}]
One more word for saying what happens if.
\end{grammarlens}

\subsection*{Your turn}
\begin{itemize}
\raggedright
  \item \emph{Say it:} \emph{ṭikānā}
  \item \emph{Say it:} \emph{ṭikānā}, once more
  \item \emph{Say it:} \emph{Usne sīṛhī dīvār se ṭikā dī.} — He leaned the ladder against the wall.
  \item \emph{Recall:} say the Hindi for to bear fruit, then the Hindi for to sway, then say \emph{ṭikānā} again
\end{itemize}

\begin{tcolorbox}[breakable,colback=teal!4,colframe=teal!35!black,title={Before you move on}]
What does \hi{टिकाना} mean? (\textquotedblleft{}To rest (something) against; to steady\textquotedblright{}.) Say it once more.
\end{tcolorbox}
\section[\texorpdfstring{\hi{गिरावट} (girāvaṭ)}{गिरावट (girāvaṭ)}]{\hi{गिरावट} (girāvaṭ) — a fall; a drop (in prices, in level)}

\label{lesson:HI-C555-giravat}



\begin{quote}
Before the new one: say the Hindi for to sway, then the Hindi for to rest (something) against.
\end{quote}

\subsection*{You'll want to know: \hi{गिरावट}}
\textbf{\hi{गिरावट}} — \emph{girāvaṭ} — \textquotedblleft{}a fall; a drop (in prices, in level)\textquotedblright{}.

\begin{quote}
\hi{सब्ज़ियों} \hi{के} \hi{दाम} \hi{में} \hi{गिरावट} \hi{आई} \hi{है।}
\end{quote}

\begin{quote}
\emph{Sabziyoṁ ke dām meṁ girāvaṭ āī hai.}
\end{quote}

\begin{quote}
Vegetable prices have dropped.
\end{quote}

Feminine.

\begin{cousinweb}
\textbf{\hi{गिरावट}} (\emph{girāvaṭ}) is the noun of \emph{girnā}, \textquotedblleft{}to fall\textquotedblright{}, which comes from \emph{*girati}, \textquotedblleft{}drips, falls\textquotedblright{}, a form CDIAL reconstructs (CDIAL 4159), with the ending -āvaṭ that makes nouns of verbs.
\end{cousinweb}

\begin{grammarlens}[title={What you've built}]
One more word for saying what happens if.
\end{grammarlens}

\subsection*{Your turn}
\begin{itemize}
\raggedright
  \item \emph{Say it:} \emph{girāvaṭ}
  \item \emph{Say it:} \emph{girāvaṭ}, once more
  \item \emph{Say it:} \emph{Sabziyoṁ ke dām meṁ girāvaṭ āī hai.} — Vegetable prices have dropped.
  \item \emph{Recall:} say the Hindi for to sway, then the Hindi for to rest (something) against, then say \emph{girāvaṭ} again
\end{itemize}

\begin{tcolorbox}[breakable,colback=teal!4,colframe=teal!35!black,title={Before you move on}]
What does \hi{गिरावट} mean? (\textquotedblleft{}A fall; a drop (in prices, in level)\textquotedblright{}.) Say it once more.
\end{tcolorbox}
\section[\texorpdfstring{\hi{पतन} (patan)}{पतन (patan)}]{\hi{पतन} (patan) — a downfall; a decline}

\label{lesson:HI-C555-patan}



\begin{quote}
Before the new one: say the Hindi for to rest (something) against, then the Hindi for a fall.
\end{quote}

\subsection*{You'll want to know: \hi{पतन}}
\textbf{\hi{पतन}} — \emph{patan} — \textquotedblleft{}a downfall; a decline\textquotedblright{}.

\begin{quote}
\hi{बुरी} \hi{आदतें} \hi{इंसान} \hi{के} \hi{पतन} \hi{का} \hi{कारण} \hi{बनती} \hi{हैं।}
\end{quote}

\begin{quote}
\emph{Burī ādteṁ insān ke patan kā kāraṇ bantī haiṁ.}
\end{quote}

\begin{quote}
Bad habits bring about a person's downfall.
\end{quote}

Masculine.

\begin{cousinweb}
\textbf{\hi{पतन}} (\emph{patan}) is a tatsama, taken whole from Sanskrit \emph{patana}, \textquotedblleft{}fall, decline, ruin\textquotedblright{} (Monier-Williams).
\end{cousinweb}

\begin{grammarlens}[title={What you've built}]
One more word for saying what happens if.
\end{grammarlens}

\subsection*{Your turn}
\begin{itemize}
\raggedright
  \item \emph{Say it:} \emph{patan}
  \item \emph{Say it:} \emph{patan}, once more
  \item \emph{Say it:} \emph{Burī ādteṁ insān ke patan kā kāraṇ bantī haiṁ.} — Bad habits bring about a person's downfall.
  \item \emph{Recall:} say the Hindi for to rest (something) against, then the Hindi for a fall, then say \emph{patan} again
\end{itemize}

\begin{tcolorbox}[breakable,colback=teal!4,colframe=teal!35!black,title={Before you move on}]
What does \hi{पतन} mean? (\textquotedblleft{}A downfall; a decline\textquotedblright{}.) Say it once more.
\end{tcolorbox}
